% 第 4 章：三角函数
\chapter{三角函数}

\section{任意角与弧度制}

\begin{problem}\level{易}
把 $150^\circ$、$-210^\circ$ 化成弧度，把 $\dfrac{5\pi}{6}$ 化成角度。
\end{problem}
\begin{solution}
$150^\circ=\dfrac{5\pi}{6}$；$-210^\circ=-\dfrac{7\pi}{6}$；$\dfrac{5\pi}{6}=150^\circ$。
\end{solution}

\begin{problem}\level{易}
写出与 $-60^\circ$ 终边相同的角的集合，并找出 $-360^\circ\sim0^\circ$ 内的那个角。
\end{problem}
\begin{solution}
$\{\beta\mid\beta=-60^\circ+k\cdot360^\circ,\ k\in\mathbb{Z}\}$；在 $-360^\circ\sim0^\circ$ 内为 $-60^\circ$。
\end{solution}

\begin{problem}\level{中}
在 $0^\circ\sim360^\circ$ 内，找出与 $-400^\circ$ 终边相同的角。
\end{problem}
\begin{solution}
$-400^\circ=-360^\circ-40^\circ$，终边与 $-40^\circ$ 相同，加 $360^\circ$ 得 $320^\circ$。
\end{solution}

\begin{problem}\level{中}
把 $22^\circ30'$ 化成弧度。
\end{problem}
\begin{solution}
$22^\circ30'=22.5^\circ=22.5\times\dfrac{\pi}{180}=\dfrac{\pi}{8}$。
\end{solution}

\begin{problem}\level{中}
半径为 $6$、圆心角为 $120^\circ$ 的扇形，求弧长与面积。
\end{problem}
\begin{solution}
$120^\circ=\dfrac{2\pi}{3}$，$l=\dfrac{2\pi}{3}\times6=4\pi$；
$S=\dfrac12lr=\dfrac12\times4\pi\times6=12\pi$。
\end{solution}

\begin{problem}\level{难}
已知扇形的周长为 $10$，面积为 $6$，求其半径与圆心角（弧度）。
\end{problem}
\begin{solution}
设半径 $r$、弧长 $l$，则 $2r+l=10$，$\dfrac12lr=6$。
由 $l=10-2r$ 代入得 $\dfrac12r(10-2r)=6$，即 $r^2-5r+6=0$，解得 $r=2$ 或 $r=3$。
$r=2$ 时 $l=6$，$\alpha=\dfrac lr=3$；$r=3$ 时 $l=4$，$\alpha=\dfrac43$。
\end{solution}

\section{任意角的三角函数}

\begin{problem}\level{中}
已知角 $\alpha$ 的终边过点 $(3,4)$，求 $\sin\alpha$、$\cos\alpha$、$\tan\alpha$。
\end{problem}
\begin{solution}
$r=5$，$\sin\alpha=\dfrac45$，$\cos\alpha=\dfrac35$，$\tan\alpha=\dfrac43$。
\end{solution}

\begin{problem}\level{中}
已知角 $\alpha$ 的终边过点 $(-4,3)$，求 $\sin\alpha$、$\cos\alpha$、$\tan\alpha$。
\end{problem}
\begin{solution}
$r=5$，$\sin\alpha=\dfrac35$，$\cos\alpha=-\dfrac45$，$\tan\alpha=-\dfrac34$。
\end{solution}

\begin{problem}\level{中}
已知 $\sin\alpha=\dfrac35$，且 $\alpha$ 是第二象限角，求 $\cos\alpha$ 与 $\tan\alpha$。
\end{problem}
\begin{solution}
$\cos^2\alpha=1-\dfrac{9}{25}=\dfrac{16}{25}$，第二象限 $\cos\alpha<0$，故 $\cos\alpha=-\dfrac45$；
$\tan\alpha=-\dfrac34$。
\end{solution}

\begin{problem}\level{中}
已知 $\cos\alpha=-\dfrac{5}{13}$，且 $\alpha$ 在第三象限，求 $\sin\alpha$、$\tan\alpha$。
\end{problem}
\begin{solution}
$\sin^2\alpha=1-\dfrac{25}{169}=\dfrac{144}{169}$，第三象限 $\sin\alpha<0$，故 $\sin\alpha=-\dfrac{12}{13}$；
$\tan\alpha=\dfrac{12}{5}$。
\end{solution}

\begin{problem}\level{中}
已知 $\sin\alpha=\dfrac{12}{13}$，$\alpha$ 为第一象限角，求 $\cos\alpha$、$\tan\alpha$。
\end{problem}
\begin{solution}
$\cos\alpha=\dfrac5{13}$，$\tan\alpha=\dfrac{12}{5}$。
\end{solution}

\begin{problem}\level{中}
化简 $\sin(180^\circ+\alpha)\cos(-\alpha)$。
\end{problem}
\begin{solution}
$\sin(180^\circ+\alpha)=-\sin\alpha$，$\cos(-\alpha)=\cos\alpha$，故原式 $=-\sin\alpha\cos\alpha$。
\end{solution}

\begin{problem}\level{中}
化简 $\cos\left(\dfrac{\pi}{2}+\alpha\right)+\sin(\pi-\alpha)$。
\end{problem}
\begin{solution}
$\cos\left(\dfrac{\pi}{2}+\alpha\right)=-\sin\alpha$，$\sin(\pi-\alpha)=\sin\alpha$，故原式 $=0$。
\end{solution}

\begin{problem}\level{中}
求 $\sin\dfrac{7\pi}{6}$。
\end{problem}
\begin{solution}
$\sin\dfrac{7\pi}{6}=\sin\left(\pi+\dfrac{\pi}{6}\right)=-\sin\dfrac{\pi}{6}=-\dfrac12$。
\end{solution}

\begin{problem}\level{难}
化简 $\dfrac{\sin(2\pi-\alpha)\cos(\pi+\alpha)}{\tan(-\alpha)}$。
\end{problem}
\begin{solution}
$\sin(2\pi-\alpha)=-\sin\alpha$，$\cos(\pi+\alpha)=-\cos\alpha$，$\tan(-\alpha)=-\tan\alpha$，
故原式 $=\dfrac{(-\sin\alpha)(-\cos\alpha)}{-\tan\alpha}
=\dfrac{\sin\alpha\cos\alpha}{-\dfrac{\sin\alpha}{\cos\alpha}}=-\cos^2\alpha$。
\end{solution}

\begin{problem}\level{难}
已知 $\tan\alpha=2$，求 $\dfrac{\sin\alpha+\cos\alpha}{\sin\alpha-\cos\alpha}$。
\end{problem}
\begin{solution}
分子分母同除以 $\cos\alpha$：$\dfrac{\tan\alpha+1}{\tan\alpha-1}=\dfrac{2+1}{2-1}=3$。
\end{solution}

\section{三角函数的图象与性质}

\begin{problem}\level{中}
求 $y=\sin x$ 与 $y=\cos x$ 的单调递增区间。
\end{problem}
\begin{solution}
$y=\sin x$：$\left[-\dfrac{\pi}{2}+2k\pi,\dfrac{\pi}{2}+2k\pi\right]$（$k\in\mathbb{Z}$）；
$y=\cos x$：$[\pi+2k\pi,2\pi+2k\pi]$（$k\in\mathbb{Z}$）。
\end{solution}

\begin{problem}\level{中}
求 $y=\tan x$ 的定义域与周期。
\end{problem}
\begin{solution}
定义域 $\left\{x\mid x\neq\dfrac{\pi}{2}+k\pi,\ k\in\mathbb{Z}\right\}$；周期 $\pi$。
\end{solution}

\begin{problem}\level{中}
用五点法在区间 $[0,2\pi]$ 上画出 $y=1+2\sin x$ 的图象（写出五个关键点）。
\end{problem}
\begin{solution}
令 $x=0,\dfrac{\pi}{2},\pi,\dfrac{3\pi}{2},2\pi$，得 $y=1,3,1,-1,1$。
五个关键点为 $(0,1)$、$\left(\dfrac{\pi}{2},3\right)$、$(\pi,1)$、$\left(\dfrac{3\pi}{2},-1\right)$、$(2\pi,1)$。
\end{solution}

\begin{problem}\level{中}
判断 $y=\sin x+\cos x$ 的奇偶性，并求它的最大值。
\end{problem}
\begin{solution}
定义域 $\mathbb{R}$ 关于原点对称；$f(-x)=-\sin x+\cos x$，既不等于 $f(x)$ 也不等于 $-f(x)$，
故既不是奇函数也不是偶函数。
最大值 $=\sqrt2$（$y=\sqrt2\sin\left(x+\dfrac{\pi}{4}\right)$）。
\end{solution}

\begin{problem}\level{难}
求函数 $y=2\sin^2 x-4\sin x+1$ 在 $x\in\left[\dfrac{\pi}{6},\dfrac{5\pi}{6}\right]$ 上的最大值与最小值。
\end{problem}
\begin{solution}
令 $t=\sin x$，由 $x\in\left[\dfrac{\pi}{6},\dfrac{5\pi}{6}\right]$ 得 $t\in\left[\dfrac12,1\right]$。
$y=2t^2-4t+1=2(t-1)^2-1$，在 $\left[\dfrac12,1\right]$ 上递减。
最大值 $y\left(\dfrac12\right)=2\cdot\dfrac14-2+1=-\dfrac12$；最小值 $y(1)=-1$。
\end{solution}

\section{函数 $y=A\sin(\omega x+\varphi)$}

\begin{problem}\level{中}
求 $y=2\sin\left(2x+\dfrac{\pi}{3}\right)$ 的振幅、周期与初相。
\end{problem}
\begin{solution}
振幅 $2$，周期 $T=\dfrac{2\pi}{2}=\pi$，初相 $\dfrac{\pi}{3}$。
\end{solution}

\begin{problem}\level{中}
求 $y=4\sin\left(3x-\dfrac{\pi}{6}\right)$ 的振幅、周期与初相。
\end{problem}
\begin{solution}
振幅 $4$，周期 $T=\dfrac{2\pi}{3}$，初相 $-\dfrac{\pi}{6}$。
\end{solution}

\begin{problem}\level{中}
求 $y=3\sin\left(\dfrac{x}{2}-\dfrac{\pi}{4}\right)$ 的振幅、周期、初相，并说明变换过程。
\end{problem}
\begin{solution}
振幅 $3$，周期 $T=4\pi$，初相 $-\dfrac{\pi}{4}$。
$y=\sin x$ 右移 $\dfrac{\pi}{4}$，横坐标伸长为原来的 $2$ 倍，纵坐标伸长为原来的 $3$ 倍。
\end{solution}

\begin{problem}\level{难}
函数 $y=A\sin(\omega x+\varphi)$ 的图象相邻两条最高点距离为 $\pi$，且过点 $\left(0,\dfrac12\right)$，
振幅为 $1$，写出它的一个解析式。
\end{problem}
\begin{solution}
周期 $T=\pi$，故 $\omega=\dfrac{2\pi}{T}=2$。振幅 $A=1$。
过 $\left(0,\dfrac12\right)$：$\sin\varphi=\dfrac12$，取 $\varphi=\dfrac{\pi}{6}$。
一个解析式为 $y=\sin\left(2x+\dfrac{\pi}{6}\right)$。
\end{solution}

\begin{problem}\level{难}
求函数 $y=\sin x+\sqrt3\cos x$ 的最小正周期与最大值。
\end{problem}
\begin{solution}
$y=2\sin\left(x+\dfrac{\pi}{3}\right)$，最小正周期 $2\pi$，最大值 $2$。
\end{solution}

\section{三角函数模型}

\begin{problem}\level{中}
某地一天的温度可用 $T(t)=8\sin\left(\dfrac{\pi}{12}t-\dfrac{7\pi}{12}\right)+20$（$t$ 为小时）近似描述，
求这一天的最高温度、最低温度与周期。
\end{problem}
\begin{solution}
最高 $20+8=28$（℃），最低 $20-8=12$（℃），周期 $T=\dfrac{2\pi}{\pi/12}=24$（小时）。
\end{solution}

\begin{problem}\level{中}
一个弹簧振子做简谐运动，位移 $y$（cm）与时间 $t$（s）满足 $y=5\sin(2\pi t)$。
求振幅、周期，并求 $t=0.25$ s 时的位移。
\end{problem}
\begin{solution}
振幅 $5$ cm，周期 $1$ s；$t=0.25$ 时 $y=5\sin\dfrac{\pi}{2}=5$ cm。
\end{solution}

\begin{problem}\level{难}
某港口一天的水深 $h$（m）与时间 $t$（h）近似满足 $h=6+4\sin\left(\dfrac{\pi}{6}t\right)$。
求一天中水深的最小值与最大值，并求水深大于 $8$ m 的时间段。
\end{problem}
\begin{solution}
最大值 $6+4=10$ m，最小值 $6-4=2$ m。
$h>8\iff\sin\left(\dfrac{\pi}{6}t\right)>\dfrac12
\iff\dfrac{\pi}{6}+2k\pi<\dfrac{\pi}{6}t<\dfrac{5\pi}{6}+2k\pi$
$\iff1+12k<t<5+12k$（$k\in\mathbb{Z}$）。
故一天中水深大于 $8$ m 的时间段为 $(1,5)$ 与 $(13,17)$ 与 $(25,29)\dots$，每次持续 $4$ 小时。
\end{solution}
